\documentclass[11pt, a4paper]{article}

\usepackage[margin=1in]{geometry}
\usepackage{graphicx}
\usepackage{booktabs}
\usepackage{hyperref}
\usepackage{titlesec}
\usepackage{xcolor}
\usepackage{amsmath}
\usepackage{enumitem}

% Formatting
\hypersetup{
    colorlinks=true,
    linkcolor=blue,
    filecolor=magenta,      
    urlcolor=cyan,
}
\titleformat{\section}{\large\bfseries}{\thesection}{1em}{}
\titleformat{\subsection}{\normalsize\bfseries}{\thesubsection}{1em}{}

% Compact spacing
\setlength{\parskip}{0.4em}
\setlength{\parindent}{0pt}
\titlespacing*{\section}{0pt}{1em}{0.4em}
\titlespacing*{\subsection}{0pt}{0.8em}{0.3em}

\begin{document}

% --- COVER PAGE ---
\begin{titlepage}
    \centering
    \vspace*{5cm}
    
    {\Huge \textbf{Predicting Port Congestion and \\ Shipment Delays Using AI}\par}
    \vspace{1.5cm}

    {\LARGE \textbf{Maritime Data Solutions}\par}
    \vspace{1.5cm}
    
    {\Large \textbf{Predictive Analytics \& Machine Learning \\ Industry Case Study}\par}
    \vspace{2cm}
    
    {\Large \textbf{Prepared by:} Nguyen Duc Tuan Dat\par}
    \vspace{0.5cm}
    
    {\large \textbf{Date:} April 2026\par}
    
    \vfill
\end{titlepage}

% --- TABLE OF CONTENTS ---
\tableofcontents
\newpage

% --- MAIN REPORT (3 pages of content) ---

\textit{\small Note: This report summarizes key design decisions and results. For detailed step-by-step analysis---including all visualizations, SHAP charts, EDA plots, and inline commentary---please refer to the fully annotated Jupyter notebook (\texttt{notebooks/01\_comprehensive\_analysis.ipynb}).}

\section{Overview}

Inspired by the severe bottlenecks at Cat Lai Port (Ho Chi Minh City) \cite{vietnamnet_catlai}, I engineered a realistic synthetic dataset modeling non-linear maritime physics---exponential yard capacity saturation \cite{mm1_queue} and weather-induced crane lockouts \cite{pema_wind}---then benchmarked three classifiers to satisfy three requirements:
\begin{enumerate}[nosep, leftmargin=1.5em]
    \item \textbf{Predict the likelihood of congestion} (delays $> 240$ min).
    \item \textbf{Provide a confidence score} via probability calibration (\texttt{predict\_proba}).
    \item \textbf{Explain key contributing factors} using SHAP, ensuring transparent, actionable output.
\end{enumerate}

\section{Approach}

\subsection{Data Generation Strategy}

To avoid the mathematical noise inherent in raw AIS tracking data, I constructed a controlled synthetic pipeline to rigorously benchmark model limitations. The generation was parameterized directly on Cat Lai Port's operational realities and \textbf{evolved across two versions}:

\begin{itemize}[nosep, leftmargin=1.5em]
    \item \textbf{V1 (Linear Baseline):} My initial dataset used a simple additive formula: $Delay = \beta_0 + \text{Capacity} \cdot w_1 + \text{Wind} \cdot w_2 + \epsilon$. This intentionally simplified model served as a sanity check---if a basic algorithm could not solve V1, my pipeline had fundamental errors. All models passed easily, which confirmed the pipeline was sound but also revealed a critical insight: V1 was \textit{too simple} to stress-test algorithms for real-world deployment.
    
    \item \textbf{V2 (Cat Lai Non-Linear Physics):} Motivated by that insight, I researched actual Cat Lai operations \cite{saigon_newport} and redesigned the generator to reflect real bottleneck mechanics. V2 integrates queueing theory with exponential capacity saturation (delays compound cubically with yard density) \cite{mm1_queue}, environmental ``hard-stops'' ($>$35-knot crane lockouts) \cite{pema_wind}, and conditionally compounding interactions---for example, severe ``midnight gridlock'' penalties \cite{vietnamnet_catlai} injected when a ship arrives at night while the port exceeds 80\% capacity. This evolution from V1 to V2 is the backbone of the project: \textit{the data complexity drives the model selection}, not the other way around.
\end{itemize}

\textbf{Ground Truth:} ``Congestion'' is defined strictly as delay $>$ 240 minutes, mapping to standard demurrage triggers and missed tidal-departure windows at this river port \cite{saigon_newport}.

\subsection{Model Selection \& Evaluation}

Because port operations experience class imbalance, Accuracy alone is deceptive. I evaluated all models using \textbf{ROC-AUC}, which rigorously measures the ability to output a calibrated confidence score across all risk thresholds.

My model selection evolved \textit{iteratively alongside the data complexity}---not as a predetermined choice:

\begin{itemize}[nosep, leftmargin=1.5em]
    \item \textbf{Step 1 --- Logistic Regression on V1:} I began with the simplest viable classifier. Logistic Regression achieved $>$98\% ROC-AUC on the linear V1 data, confirming that the pipeline and features were correctly engineered. However, subsequent research into Cat Lai's actual physics \cite{saigon_newport, vietnamnet_catlai} revealed that V1 was structurally insufficient---real-world logistics bottlenecks do not scale linearly \cite{mm1_queue}. A model solving a linear simulation would immediately fail in a non-linear port environment.
    
    \item \textbf{Step 2 --- The V2 Pivot:} This realization prompted the creation of the V2 dataset and forced a model paradigm shift. I introduced \textbf{Random Forest} as a mid-complexity candidate. Its independent decision trees effectively navigated ``hard-stop'' threshold rules (like 35-knot crane lockouts \cite{pema_wind}, captured via \texttt{max\_depth=5} trees). However, because each tree votes independently, Random Forest struggled with the \textit{conditionally compounding} multi-variable interactions in V2---specifically the ``midnight gridlock'' effect \cite{vietnamnet_catlai} where nighttime arrival \textit{multiplies} (not merely adds to) high-capacity delays.
    
    \item \textbf{Step 3 --- XGBoost (Champion):} Drawing on gradient boosting literature, I implemented \textbf{XGBoost} with \texttt{binary:logistic} objective. Unlike Random Forest's parallel trees, XGBoost builds \textit{sequentially}---each new tree specifically corrects the errors of the previous ensemble. This error-correcting architecture proved critical for V2: XGBoost dominated all trials, pushing beyond \textbf{99\% ROC-AUC}, because it could learn precisely how nighttime schedules \textit{exponentially multiply} the severity of high capacity---the exact interaction that defeated simpler models.
\end{itemize}

\subsection{Interpretability Strategy}

An opaque score of ``98\% Congested'' provides zero mitigation guidance. I integrated \textbf{SHAP (SHapley Additive exPlanations)} at two tiers to ensure predictions are not only accurate but operationally actionable:
\begin{itemize}[nosep, leftmargin=1.5em]
    \item \textbf{Global (Strategic):} Beeswarm summaries confirmed XGBoost correctly prioritized \texttt{queue\_length} and \texttt{wind\_speed\_knots} as dominant drivers---validating it independently recovered Cat Lai's physics from raw data alone.
    \item \textbf{Local (Tactical):} Waterfall charts decompose individual predictions feature-by-feature, translating each score into a specific, actionable recommendation (e.g., \textit{``halt berthing until wind drops below the 35-knot safety threshold''} \cite{pema_wind}).
\end{itemize}

\section{Assumptions}

\textbf{Definition of Congestion:} Any delay $>$ \textbf{240 minutes}. Cat Lai is a river port on the Dong Nai with semi-diurnal tidal constraints; a 4-hour delay realistically causes a vessel to miss its narrow tidal departure window \cite{saigon_newport}.

\textbf{Synthetic Data Physics:} All V2 non-linear rules are grounded in documented realities:
\begin{itemize}[nosep, leftmargin=1.5em]
    \item \textbf{Night Gridlock:} Cat Lai processes $>$22,000 vehicles/day; HCMC restricts heavy trucks during peak hours, concentrating traffic into off-peak windows \cite{vietnamnet_catlai}.
    \item \textbf{Exponential Capacity Saturation:} Based on $M/M/1$ queueing theory ($W = \rho / (1-\rho)$); modeled as a cubic multiplier ($1 + \rho^3$) on base processing time---creating exponential delay growth as yard density increases. Beyond 80\%, this compounds with the ``midnight gridlock'' interaction \cite{mm1_queue}.
    \item \textbf{Environmental Hard-Stop ($>$35 knots):} Derived from industry crane safety standards; our 35-knot ($\sim$18~m/s) threshold is a conservative operational limit \cite{pema_wind}.
\end{itemize}

\section{Results}

\subsection{Model Performance \& Scalability}

Table~\ref{tab:roc_auc} presents ROC-AUC across both dataset versions, scaled from 500 to 500,000 records. Two critical patterns validate the iterative design described above.

\begin{table}[h!]
\centering
\caption{ROC-AUC (\%) Across Dataset Versions and Sizes}
\label{tab:roc_auc}
\small
\begin{tabular}{llccc}
\toprule
\textbf{Version} & \textbf{Records} & \textbf{Logistic Reg.} & \textbf{Random Forest} & \textbf{XGBoost} \\
\midrule
V1 -- Linear    & 500       & 95.62 & 92.66 & 93.79 \\
               & 5,000     & 98.81 & 95.63 & 98.78 \\
               & 50,000    & 98.49 & 94.19 & 99.36 \\
               & 500,000   & 98.60 & 93.98 & 99.60 \\
\midrule
V2 -- Non-Linear & 500      & 95.32 & 94.68 & 96.05 \\
                & 5,000    & 94.71 & 98.37 & 99.51 \\
                & 50,000   & 94.77 & 98.13 & 99.69 \\
                & 500,000  & 94.74 & 98.14 & \textbf{99.76} \\
\bottomrule
\end{tabular}
\end{table}

\textbf{Pattern 1 (V1):} All models succeed ($>$92\%) on linear data---the simple additive relationships are learnable by any classifier, confirming the V1 sanity check.

\textbf{Pattern 2 (V2):} Logistic Regression \textbf{plateaus at $\sim$94.7\%} regardless of scale---it structurally cannot capture the exponential interactions modeled by queueing theory \cite{mm1_queue}. XGBoost \textbf{scales from 96.05\% to 99.76\%}, with performance improving monotonically as data grows. This $\sim$5-point gap represents the difference between a model that can reverse-engineer real port physics and one that cannot. XGBoost is selected as the champion model.

\subsection{Explainability and Actionable Insights}

SHAP decomposes every prediction into feature-level contributions, delivering two tiers of operational value:

\textbf{Global (Beeswarm):} Across the V2 test set, \texttt{queue\_length} and \texttt{wind\_speed\_knots} emerged as the dominant congestion drivers---confirming XGBoost independently discovered Cat Lai's capacity saturation \cite{mm1_queue} and weather threshold mechanics \cite{pema_wind} without explicit rules.

\textbf{Local (Waterfall -- Vessel \#0):} For a vessel classified ``Clear'' (0.5\% congestion confidence), SHAP decomposes the prediction from baseline $E[f(X)] = 2.079$ to final $f(x) = -5.254$. The dominant risk-reducing driver was \texttt{arrival\_vessels\_24h = 13} (low vessel density in the 24-hour window), with favorable wind and moderate queue lengths further pushing the score toward ``Clear.'' This transforms a black-box number into a \textit{feature-by-feature narrative}---port operators see not just \textit{what} the model predicts, but \textit{why}, and which specific operational levers to adjust.

Together, these two layers deliver \textbf{macro-level strategic visibility} (system-wide drivers) and \textbf{micro-level tactical guidance} (per-vessel explanations with actionable recommendations).

\section{Limitations \& Future Work}

Three honest limitations exist:
\begin{itemize}[nosep, leftmargin=1.5em]
    \item \textbf{Synthetic Data Ceiling:} The 99\%+ ROC-AUC reflects evaluation against the same physics rules used for training. Real AIS feeds introduce sensor noise, missing values, and unmodeled human behavior---expect realistic production performance in the 75--85\% range.
    \item \textbf{Single Split:} Results use one 80/20 split with a fixed random seed, not cross-validated.
    \item \textbf{No Temporal Memory:} Each arrival is treated independently; real congestion cascades over time.
\end{itemize}

\textbf{Given more time}, I would improve in two areas: (1)~\textit{Data Engineering}---introduce controlled noise, missing values, and lag features (e.g., average delay over the past 3~hours) into the synthetic pipeline to train a more resilient model; (2)~\textit{Model Training}---replace fixed hyperparameters with Bayesian Optimization or GridSearchCV and implement 5-fold cross-validation for statistically robust performance estimates.

% --- REFERENCES (separate page) ---
\newpage
\section{References}
The V2 synthetic constraints were parameterized using documented maritime physics and local HCMC traffic regulations:

\begin{thebibliography}{9}
\bibitem{vietnamnet_catlai}
VietNamNet Global, \textit{``Vietnam's largest port faces serious road congestion,''} Oct.\ 2023. Confirms that Cat Lai Port handles ``over 70\% of containers for international trade'' with ``some 22,000 vehicles transporting containers to the port every day,'' and documents multi-hour traffic jams on approach roads lasting over 5~hours. \\
\url{https://vietnamnet.vn/en/vietnam-s-largest-port-faces-serious-road-congestion-2202239.html}

\bibitem{mm1_queue}
$M/M/1$ Queueing Model --- Classical Operations Research. In the $M/M/1$ framework, mean wait time degrades as $W_q = \rho / \mu(1 - \rho)$: at 80\% server utilization, the mean wait is 4$\times$ the service time; at 90\%, it is 9$\times$. This exponential degradation curve is the mathematical basis for our capacity saturation threshold in the V2 synthetic generator. \textit{See: Section ``Average number of customers in the system,'' Response Time formula.} \\
\url{https://en.wikipedia.org/wiki/M/M/1_queue}

\bibitem{pema_wind}
TT Club / PEMA / ICHCA International, \textit{``Port \& Terminal Safety --- Loss Prevention Guidance.''} Joint industry guidance recommending a maximum operating wind speed of $\sim$22~m/s ($\approx$43~knots) for quay container cranes, with storm-securing procedures at 35~m/s. Our 35-knot ($\sim$18~m/s) threshold is a conservative limit below the operational maximum. \\
\url{https://www.ttclub.com/loss-prevention/port-safety/}

\bibitem{saigon_newport}
Vietnam Ports Association (VPA), \textit{``Port of Ho Chi Minh City --- Cat Lai Terminal.''} Official port directory confirming Cat Lai's classification as a river port on the Dong Nai River with berth depth of $\sim$12--12.5m, irregular semi-diurnal tidal regime (2.5--4m range), and draft-constrained navigation requiring tidal-window scheduling. \\
\url{https://www.vpa.org.vn}
\end{thebibliography}

\end{document}
